\documentclass[11pt]{article}
\usepackage[margin=1in]{geometry}
\usepackage{amsmath,amssymb,amsthm}
\usepackage{enumitem}
\usepackage[hidelinks]{hyperref}
\usepackage{mathtools}

\newtheorem{theorem}{Theorem}
\newtheorem{lemma}[theorem]{Lemma}
\newtheorem{proposition}[theorem]{Proposition}
\newtheorem{corollary}[theorem]{Corollary}
\newtheorem{definition}[theorem]{Definition}
\newtheorem{remark}[theorem]{Remark}

%% -- Notation --
\newcommand{\pair}[2]{\langle #1, #2 \rangle}
\newcommand{\pmark}{\pair{O}{O}}   % the "one" / false marker
\newcommand{\mark}{\mathbf{1}}
\newcommand{\SchemaF}{\textbf{Schema~F}}
\newcommand{\embed}{\ulcorner}
\newcommand{\embedr}{\urcorner}
\newcommand{\code}[1]{\embed #1 \embedr}
\newcommand{\geval}{\mathrm{geval}}
\newcommand{\thFun}{\mathrm{thFun}}
\newcommand{\nat}[1]{\overline{#1}}

\title{The Boundary of Nelson's Feasibility Program:\\
  Per-Program Proofs Are Feasible,\\
  Self-Reference Is Not}

\author{}
\date{}

\begin{document}
\maketitle

\begin{abstract}
We formalize in Agda a tree-based equational system combining
Guard's syntax (binary trees with functors) and Rose's
purely equational proof style~\cite{rose1953} (no propositional
connectives --- only equations and tree induction).
The formalization contains:
(1)~a diagonal G\"odel~I proof~\cite{godel1931}
(the G\"odel sentence is not provable; 168 lines, 0 postulates);
(2)~a micro-Nelson instance~\cite{nelson1986} verifying 56
programs against a fixed target;
(3)~a uniform proof-construction mechanism (the R/Q technique)
producing constant-size, BLevel-0 invariant proofs for all
oblivious step functions;
(4)~the Kritchman--Raz contradiction theorem~\cite{kritchmanraz2010},
fully proved.
The main conclusion: the \emph{per-program} content of a consistency
proof is feasible ($O(1)$, ground-verifiable), but combining these
proofs into a system-level consistency statement requires
self-reference, which is irreducibly infeasible.
\end{abstract}

%% ---------------------------------------------------------------
\section{The formal system}

The system combines two ideas: \emph{Guard}'s syntax of binary trees
with typed functors, and \emph{Rose}'s Goodstein-style proof
theory~\cite{rose1953} where the only judgement form is an equation
$t = u$ between terms.  There are no propositional connectives,
no sequents, no logical rules --- only equations and a tree-induction
principle (\SchemaF).

%% -- Trees --
\subsection*{Trees}

The ground data are finite binary trees:
\[
  a, b \;::=\; O \;\mid\; \pair{a}{b}
\]
where $O$ is the leaf and $\pair{a}{b}$ is the node.
We write $\nat{n}$ for the unary encoding of $n \in \mathbb{N}$:
$\nat{0} = O$, $\nat{n+1} = \pair{O}{\nat{n}}$.

%% -- Terms and functors --
\subsection*{Terms and functors}

Terms are built from trees, variables, and functor applications.
We use $t, u, v, \ldots$ for terms and $x, y$ for variables.
There are two sorts of functors: unary ($f, g : \mathrm{Fun}_1$)
and binary ($h : \mathrm{Fun}_2$).

\[
\begin{array}{rcl@{\qquad}l}
  f, g &::=& \mathrm{I} & \text{identity} \\
       &\mid& \pi_1 & \text{first projection} \\
       &\mid& \pi_2 & \text{second projection} \\
       &\mid& f \circ g & \text{composition} \\
       &\mid& h[f, g] & \text{binary composition: } t \mapsto h(f(t), g(t)) \\
       &\mid& \mathrm{R}(z, s) & \text{tree recursion with base } z \text{ and step } s \\
       &\mid& \mathrm{K}_c & \text{constant } c \\[6pt]
  h &::=& \mathrm{P} & \text{pairing: } (t, u) \mapsto \pair{t}{u} \\
    &\mid& \mathrm{C} & \text{left projection: } (t, u) \mapsto t \\
    &\mid& \mathrm{Lf}(f) & \text{lift: } (t, u) \mapsto f(t) \\
    &\mid& \mathrm{Ps}(f, h) & \text{post-compose: } (t, u) \mapsto f(h(t, u)) \\
    &\mid& \mathrm{Fn}(h_1, h_2, h) & \text{fan: } (t, u) \mapsto h(h_1(t,u),\, h_2(t,u)) \\
    &\mid& \mathrm{If} & \text{conditional} \\
    &\mid& \mathrm{Eq} & \text{tree equality test}
\end{array}
\]

Terms:
\[
  t, u \;::=\; O \;\mid\; x \;\mid\; f(t) \;\mid\; h(t, u)
\]

An \emph{equation} is a pair $t = u$ of terms.

%% -- Embedding --
\subsection*{Embedding trees into terms}

Every tree $a$ determines a \emph{closed term} $\hat{a}$:
\[
  \hat{O} = O, \qquad \widehat{\pair{a}{b}} = \pair{\hat a}{\hat b}.
\]
(In the Agda code this function is called \texttt{reify}.)
Conversely, every closed term reduces to $\hat{a}$ for a unique
tree~$a$. The embedding $a \mapsto \hat{a}$ lets us use trees as
term-level constants.

%% ---------------------------------------------------------------
\subsection*{Derivation rules}

A derivation is a well-founded tree built from the rules below.
Every rule may use a hypothesis equation $H$;
we write $H \vdash t = u$ for ``$t = u$ is derivable from~$H$.''

\medskip\noindent
\textbf{Computation axioms.}
\begin{alignat*}{2}
  &\text{Identity:}          &\quad \mathrm{I}(t) &= t \\
  &\text{First projection:}  &\quad \pi_1(\pair{t}{u}) &= t \\
  &\text{Second projection:} &\quad \pi_2(\pair{t}{u}) &= u \\
  &\text{Constant:}          &\quad \mathrm{C}(t, u) &= t \\
  &\text{Composition:}       &\quad (f \circ g)(t) &= f(g(t)) \\
  &\text{Binary composition:}&\quad h[f,g](t) &= h(f(t),\, g(t)) \\
  &\text{Lift:}              &\quad \mathrm{Lf}(f)(t, u) &= f(t) \\
  &\text{Post-compose:}      &\quad \mathrm{Ps}(f, h)(t, u) &= f(h(t,u)) \\
  &\text{Fan:}               &\quad \mathrm{Fn}(h_1,h_2,h)(t,u)
                                      &= h(h_1(t,u),\, h_2(t,u)) \\
  &\text{Constant functor:}  &\quad \mathrm{K}_c(t) &= c
\end{alignat*}

\noindent
\textbf{Tree recursion.}
\begin{alignat*}{2}
  &\text{Base:} &\quad \mathrm{R}(z, s)(O) &= z \\
  &\text{Step:} &\quad \mathrm{R}(z, s)(\pair{t}{u})
    &= s\bigl(\pair{t}{u},\;
         \pair{\mathrm{R}(z,s)(t)}{\mathrm{R}(z,s)(u)}\bigr)
\end{alignat*}

\noindent
\textbf{Conditional.}
\begin{alignat*}{2}
  &\text{Leaf branch:}     &\quad \mathrm{If}(O,\, \pair{t}{u}) &= t \\
  &\text{Node branch:}     &\quad \mathrm{If}(\pair{t_1}{t_2},\, \pair{u_1}{u_2}) &= u_2
\end{alignat*}

\noindent
\textbf{Tree equality.}
\begin{alignat*}{2}
  &\mathrm{Eq}(O, O) &= O \\
  &\mathrm{Eq}(O, \pair{t}{u}) &= \pmark \\
  &\mathrm{Eq}(\pair{t}{u}, O) &= \pmark \\
  &\mathrm{Eq}(\pair{t_1}{t_2}, \pair{u_1}{u_2})
    &= \mathrm{If}\bigl(\mathrm{Eq}(t_1, u_1),\;
         \pair{\mathrm{Eq}(t_2, u_2)}{\pmark}\bigr)
\end{alignat*}
Here $\mathrm{Eq}$ returns $O$ for equal trees and $\pmark$ for unequal.

\medskip\noindent
\textbf{Structural rules.}
\begin{gather*}
  \frac{\phantom{t=t}}{t = t}\;\text{(refl)}
  \qquad\qquad
  \frac{t = u}{u = t}\;\text{(sym)}
  \qquad\qquad
  \frac{t = u \qquad u = v}{t = v}\;\text{(trans)}
  \\[8pt]
  \frac{t = u}{f(t) = f(u)}\;\text{(cong$_1$)}
  \qquad\qquad
  \frac{t = u}{h(t, v) = h(u, v)}\;\text{(cong$_L$)}
  \qquad\qquad
  \frac{t = u}{h(v, t) = h(v, u)}\;\text{(cong$_R$)}
\end{gather*}

\noindent
\textbf{Substitution.}
\[
  \frac{H \vdash t = u}
       {H \vdash t[c/x] = u[c/x]}
  \;\text{(inst)}
\]

\noindent
\textbf{Hypothesis.}
\[
  \frac{\phantom{H}}{H \vdash H}
  \;\text{(hyp)}
\]

\noindent
\textbf{Schema~F} (tree induction).
If $f$ and $g$ are unary functors, $z$ a term, $s$ a binary functor,
and we have derivations of:
\begin{align*}
  f(O) &= z,  &  f(\pair{x}{y}) &= s(\pair{x}{y},\; \pair{f(x)}{f(y)}), \\
  g(O) &= z,  &  g(\pair{x}{y}) &= s(\pair{x}{y},\; \pair{g(x)}{g(y)}),
\end{align*}
then $f(x) = g(x)$.

\medskip
\SchemaF{} expresses the uniqueness of tree recursion:
if $f$ and $g$ satisfy the same recurrence equations
(same base~$z$, same step~$s$), they agree on all inputs.
Restricted to the unary trees $\nat{n}$, this gives ordinary
mathematical induction.

\subsection*{Consistency}

The system is \emph{consistent} if there is no derivation of
$O = \pmark$.  Since $O$ and $\pmark$ are distinct closed terms
(leaf vs.\ node), consistency means the system does not confuse
structurally different trees.

%% ---------------------------------------------------------------
\section{Proof encoding and ground evaluation}

Every derivation rule is assigned a \emph{tag} (a natural number
$0$--$25$).  A tree $a$ is interpreted as a proof encoding:
the function $\thFun : \mathrm{Tree} \to \mathrm{Tree}$ reads the
tag from~$a$ and reconstructs the equation that~$a$ purports to
prove, producing $\pair{\code{t}}{\code{u}}$ where $\code{t}$
is the tree representation of the term~$t$.

The ground evaluator $\geval$ takes a coded term (a tree representing
a term built from functor applications) and computes its value by
dispatching on functor-code tags (numbered $26$--$32$,
disjoint from the proof tags).  For \emph{ground} coded terms
(containing no variables), $\geval$ correctly computes the semantic
value.  For terms with variables, $\geval$ returns a stuck
(passthrough) result.

The proof checker $\mathrm{checkProof}_2(a)$ runs $\thFun$ to extract
the equation encoded by~$a$, then checks whether $\geval$ produces
the same tree on both sides.  The distinction between ground and
stuck is exactly the distinction between \emph{oblivious} and
\emph{non-oblivious} proofs.

%% ---------------------------------------------------------------
\section{Programs and iteration}

A \emph{program} is a unary functor of the form
$\mathrm{iterate}(f, O) = \mathrm{R}(O, s_f)$
where $f$ is a step function and $s_f$ extracts the right recursive
result and applies~$f$.  It satisfies:
\begin{align*}
  \mathrm{iterate}(f, O)(O) &= O, \\
  \mathrm{iterate}(f, O)(\pair{t}{u}) &= f\bigl(\mathrm{iterate}(f, O)(u)\bigr).
\end{align*}
Applied to $\nat{n}$, this gives $f^n(O)$.
The \emph{code size} of a program~$p$ is $|p|$, the number of
nodes in the tree encoding of $p$'s functor code.

%% ---------------------------------------------------------------
\section{The R/Q technique}

Two lemmas, both proved by \SchemaF{} at BLevel~$0$
(no variable instantiation):

\begin{lemma}[R --- collapse]
  For all terms $t$:
  \[
    \mathrm{If}(t,\; \pair{\pmark}{\pmark}) \;=\; \pmark.
  \]
\end{lemma}
\begin{proof}
  Define $R(t) = \mathrm{If}(t, \pair{\pmark}{\pmark})$.
  Base: $R(O) = \pmark$ (leaf branch of~$\mathrm{If}$).
  Step: $R(\pair{x}{y}) = \pmark$ (node branch).
  Both match $\mathrm{K}_{\pmark}$.
  \SchemaF{} gives $R = \mathrm{K}_{\pmark}$.
\end{proof}

\begin{lemma}[Q --- for $a \ne b$ with one component $O$]
  Let $a = \pmark$, $b = O$.  For all terms~$t$:
  \[
    \mathrm{If}\!\bigl(\mathrm{Eq}(t, a),\;
      \pair{\mathrm{Eq}(t, b)}{\pmark}\bigr) \;=\; \pmark.
  \]
\end{lemma}
\begin{proof}
  Define $Q(t) = \mathrm{If}(\mathrm{Eq}(t, \pmark),\;
    \pair{\mathrm{Eq}(t, O)}{\pmark})$.
  Base: $Q(O)$. Since $\mathrm{Eq}(O, \pmark) = \pmark$,
  $\mathrm{If}$ takes the node branch, giving~$\pmark$.
  Step: $Q(\pair{x}{y})$. Since $\mathrm{Eq}(\pair{x}{y}, O) = \pmark$,
  the pair argument becomes $\pair{\pmark}{\pmark}$.
  By~R, the result is~$\pmark$.
  \SchemaF{} closes: $Q = \mathrm{K}_{\pmark}$.
\end{proof}

The key: $Q$ captures a case analysis without case-splitting.
When $t = a$, the first comparison returns~$O$ and
$\mathrm{If}$ falls through to the second comparison, which detects
$t \ne b$ (since $a \ne b$).  When $t \ne a$, the first comparison
returns~$\pmark$ and $\mathrm{If}$ short-circuits.
Both branches produce~$\pmark$, so \SchemaF{} applies.

%% ---------------------------------------------------------------
\section{Three proof families}

For step functions $f = \mathrm{P}[g, h]$ where
$g, h \in \{I, \pi_1, \pi_2, \mathrm{K}_O, \mathrm{K}_{\pmark}\}$,
and target $\xi = \pair{\pmark}{O}$,
three parametric proof constructors handle all oblivious cases.
In each case, the theorem proved is:
for all terms~$t$,
$\mathrm{Eq}(\mathrm{iterate}(f, O)(t),\; \hat\xi) = \pmark$.

\subsection*{Family A: Constant $g$, first-component mismatch}
If $g = \mathrm{K}_c$ with $c \ne \pmark$
(so $\mathrm{Eq}(c, \pmark) = \pmark$),
then in the step case:
\[
  \mathrm{Eq}\bigl(\pair{c}{\ldots},\; \pair{\pmark}{O}\bigr)
  \;=\; \mathrm{If}\!\bigl(\underbrace{\mathrm{Eq}(c, \pmark)}_{=\;\pmark},\;
    \ldots\bigr)
  \;=\; \pmark.
\]
The first-component comparison short-circuits.
\textbf{5 programs} (all with $g = \mathrm{K}_O$).

\subsection*{Family B: Non-constant $g$, $h = \mathrm{K}_{\pmark}$}
The second-component comparison gives
$\mathrm{Eq}(\pmark, O) = \pmark$.
The $\mathrm{If}$ argument becomes $\pair{\pmark}{\pmark}$.
R~collapses it.
\textbf{3 programs} ($g \in \{I, \pi_1, \pi_2\}$).

\subsection*{Family C: $g = h$ (same function), Q~lemma}
$f(t) = \pair{g(t)}{g(t)}$.  The comparison with
$\xi = \pair{\pmark}{O}$ gives:
\[
  \mathrm{If}\!\bigl(\mathrm{Eq}(g(r), \pmark),\;
    \pair{\mathrm{Eq}(g(r), O)}{\pmark}\bigr)
  \;=\; Q(g(r)) \;=\; \pmark,
\]
where $r$ is the recursive result.
\textbf{3 programs} ($g = h \in \{I, \pi_1, \pi_2\}$).

\medskip\noindent
\textbf{Total:} 11 programs with full equational derivations.
All 23 $\mathrm{P}[g,h]$ combinations pass
$\mathrm{checkProof}_2$ ($\geval$~verification).

%% ---------------------------------------------------------------
\section{The micro-Nelson instance}

\texttt{MicroNelson.agda} enumerates 28 step functions
(all $\mathrm{P}[g,h]$ with
$g, h \in \{I, \pi_1, \pi_2, \mathrm{K}_O, \mathrm{K}_{\pmark}\}$,
plus 5~base functors) in two modes (iterate + direct), giving
56~programs.  The target is
$\xi_4 = \pair{\pair{\pair{O}{O}}{O}}{\pair{O}{O}}$,
chosen deeper than any $f(O)$ for depth-$1$ step functions.

All 56 satisfy
$\mathrm{checkProof}_2(\mathrm{mkTest}(p, \xi_4)) = \mathrm{true}$,
verified by Agda \texttt{refl}.
The adversarial program $\mathrm{K}_{\hat\xi_4}$
(which always outputs~$\xi_4$) correctly returns~\texttt{false}.

This is a complete pigeonhole instance at constructor-depth~$\le 1$:
every program of bounded depth is accounted for,
and $\xi_4$ is not in any program's range.

%% ---------------------------------------------------------------
\section{The Kritchman--Raz framework}

The Chaitin machine~\cite{chaitin1974} $C_l$ enumerates trees and
checks each one as a potential proof encoding.
If the system proves its own consistency, the search terminates
and $C_l$ outputs a string~$\xi$ with
$K(\xi) > l$~\cite{kolmogorov1965}.
But $|C_l|$ is bounded: the step function is fixed and $l$ enters
only as data via binary encoding, contributing $O(\log l)$ to the
code size.
For large~$l$, $|C_l| < l$, so $C_l$ itself witnesses
$K(\xi) \le |C_l| < l$, contradicting~$K(\xi) > l$.

\begin{theorem}[\texttt{krClean}, fully proved in Agda]
  Given:
  \begin{enumerate}[nosep]
    \item a Chaitin machine family
          $\mathit{cm} : \mathbb{N} \to \mathrm{Fun}_1$,
    \item a threshold $T$ with $|\mathit{cm}(l)| \le l$
          for $l \ge T$,
    \item for each $l \ge T$: a tree $\xi$, a clock $c$,
          a proof that $\mathit{cm}(l)$ outputs $\xi$ at clock~$c$,
          and a proof that no $\mathrm{Fun}_1$ of size $\le l$
          outputs~$\xi$,
  \end{enumerate}
  then $\bot$.
\end{theorem}

\noindent
The proof is direct: $\mathit{cm}(T)$ outputs $\xi$ at clock~$c$
and has $|\mathit{cm}(T)| \le T$, contradicting the
incompressibility of~$\xi$.

%% ---------------------------------------------------------------
\section{The boundary}

The formalization reveals a sharp decomposition:

\begin{center}
\begin{tabular}{lll}
\textbf{Component} & \textbf{Status} & \textbf{Obstacle} \\
\hline
Per-program invariant & Feasible & None ($O(1)$, BLevel 0) \\
Ground verification & Feasible & None ($\mathrm{checkProof}_2$) \\
Pigeonhole (finite $l$) & Feasible & None (enumeration) \\
\hline
$\forall p.\; |p| \le l \Rightarrow \ldots$ & Infeasible & Self-reference \\
Chaitin machine as $\mathrm{Fun}_1$ & Infeasible & Self-reference \\
$|C_l| < l$ for large $l$ & Infeasible & Self-reference \\
\end{tabular}
\end{center}

The transition from feasible to infeasible is precisely the
transition from ``for each specific program~$p$'' to
``for all programs of size~$\le l$.''

\begin{itemize}[nosep]
\item ``For each~$p$'': the R/Q technique produces a constant-size
  proof.  The proof is \emph{oblivious}: it examines only the
  structural mismatch between $p$'s output format and~$\xi$,
  not what $p$ computes.

\item ``For all~$p$'': requires the system to enumerate its own
  programs, count them (pigeonhole), and recognize that the
  Chaitin search procedure is itself a program (self-reference).
  This is the same obstacle as in the diagonal proof, repackaged.
\end{itemize}

The diagonal proof (\texttt{GodelII.agda}) concentrates
self-reference in the substitution machinery.
The Kritchman--Raz approach distributes it across counting and search.
Both are blocked by the same fundamental constraint:
the system cannot fully reason about its own proof structure.

%% ---------------------------------------------------------------
\section{Implications for Nelson's program}

\begin{enumerate}
\item \textbf{The per-program feasibility claim holds.}
  Each invariant proof is $O(1)$, BLevel~$0$, and ground-verifiable.
  The R/Q technique provides this uniformly for all oblivious
  programs.
  The $\geval$ evaluator extends coverage to non-oblivious
  programs at the metalevel.
  This is the positive content of Nelson's program~\cite{nelson1986}.

\item \textbf{The organizational infeasibility persists.}
  Combining per-program proofs into a system-level consistency
  statement requires self-reference --- whether via diagonalization
  or via counting.
  The Kritchman--Raz approach~\cite{kritchmanraz2010}, which avoids
  the diagonal lemma entirely, hits the same wall.
  The obstacle is not an artifact of the proof method.

\item \textbf{The boundary is intrinsic.}
  The oblivious/non-oblivious distinction is a local
  version of the each/all boundary.
  Oblivious proofs work by structure alone (the output format
  mismatches~$\xi$).
  Non-oblivious proofs require knowing the actual computed value ---
  a miniature version of the self-reference problem.
  The BLevel hierarchy is related to Aschieri and Berardi's
  backtracking levels~\cite{aschieri2011}.

\item \textbf{Bounded consistency is feasible.}
  For a fixed depth bound~$d$, the micro-Nelson instance gives
  a complete, computationally verified consistency proof:
  every program of depth~$\le d$ is checked, and all pass.
  This does not extend to unbounded~$d$ without self-reference.
\end{enumerate}

%% ---------------------------------------------------------------
\section{Formalization summary}

All results are formalized in Agda~2.9.0 with
\texttt{-{}-safe -{}-without-K -{}-exact-split}.
Zero postulates throughout.

\medskip
\begin{center}
\begin{tabular}{lrl}
\textbf{File} & \textbf{Lines} & \textbf{Content} \\
\hline
GodelII.agda & 168 & Diagonal G\"odel~I \\
MicroNelson.agda & 258 & 56 programs, $\mathrm{checkProof}_2$ = true \\
AutoProof.agda & 754 & R/Q library, 3 proof families, 11 derivations \\
Chaitin.agda & 395 & Binary encoding, tree enum., \texttt{krClean} \\
KRShortProof.agda & 310 & Case~1 template (first-component mismatch) \\
DoublerProof.agda & 430 & Case~3 template (Q~lemma) \\
GeneralQ.agda & 443 & Cases~2, 4; tripler, mixed \\
SwapProof.agda & 492 & Case~4 template ($h_2$~helper) \\
\hline
Guard/ total & ${\sim}40$ files & Full system including infrastructure \\
\end{tabular}
\end{center}

\bibliographystyle{plain}
\bibliography{references}

\end{document}
